% ============================================================================
% sec9_conclusion.tex — Paper 2: Conclusion
% ============================================================================

\section{Conclusion}
\label{sec:conclusion}

This paper develops and estimates a structural model of cover-bidder
deployment in procurement auctions. The model formalizes the cartel's
trade-off between detection risk, cover-bidding cost, and
minimum-bidder constraints, and generates five testable predictions---all
confirmed in data from 40~million bids in Brazilian public procurement
(2009--2019).

The central finding is the dispersion paradox: coordinated cover
bidding---the empirically dominant regime ($\Delta$BIC~$= -91{,}473$)---
produces \emph{lower} within-tender bid dispersion than genuine
competition ($\hat{\sigma}_c / \hat{\sigma}_g = 0.72$). This reverses
a maintained assumption in the cartel detection literature
\citep{imhof2019detecting, huber2019machine}, which associates
collusion with high dispersion. The mechanism explains why
participation-based screens (AUC~$= 0.94$) outperform
variance-based alternatives (AUC~$= 0.79$): the participation
footprint of cover bidding is invariant to the bid distribution,
while bid-level statistics are structurally uninformative under
Regime~2.

The illustrative welfare counterfactuals highlight the policy stakes.
The model's corner solution shows that the convite minimum-bidder
rule forces cover-bidder deployment in 39.2\% of FL-present
tenders; removing it would eliminate the institutional demand for
cover bidding. The welfare-maximizing detection threshold
($1.6\times$IQR) coincides with the statistical and
detection-theoretic optima, providing a coherent basis for
operational deployment.

\paragraph{Limitations.}
The reduced-form price associations are conditional correlations that
reflect both treatment and selection; we do not claim causal
identification of a causal price effect. The event-study pre-trends
confirm endogenous market selection, and the IV suffers from
balance-test violations. The structural model characterizes the
cover-bid mechanism---the distribution and deployment of cover
bids---but the price association reflects both treatment and
selection. FL persistence across time is low (10\%), consistent with
FL status capturing an environment indicator (market under cartel
influence) rather than a stable firm type. The calibrated
detection-penalty parameter $\hat{\phi}_0$ is indistinguishable
from zero, limiting quantitative counterfactuals along the
enforcement margin; this may reflect insufficient enforcement
variation in the BEC/SP data rather than true structural
insensitivity. The concentrated-NLS bootstrap achieves 58\%
convergence (291/500), adequate for standard errors but limiting
for tail inference. All results come from a single procurement
platform; replication is needed.

\paragraph{Implications for detection methodology.}
The dispersion paradox has a practical corollary: any screen that
conditions on within-tender bid statistics will underperform when
facing coordinated cartels---precisely the cartels with the highest
welfare cost. Participation-based screens, which detect the
\emph{operational footprint} of cover bidding regardless of bid
calibration, fill this gap. The structural model's specification
tests show that the IQR rule captures the most discriminative
dimension of cover bidding; adding bid-level features or structural
scores does not improve detection.

\paragraph{Future research.}
Three directions merit investigation: (i)~replication in other
procurement systems to test external validity;
(ii)~dynamic models of cartel--authority interaction, where cartels
adapt to screening by rotating cover bidders or allowing occasional
wins; and (iii)~field experiments in which procurement agencies
deploy FL screening and measure deterrence effects.
